\chapter{Keys, bindings and the network}
\label{ch:keys}

This chapter is about everything that leaves your machine on purpose: model
requests, keyed APIs your tools call, and MCP servers.

\section{The keychain}

Keys live in an encrypted file, the keychain, unlocked with a passphrase
you choose (not recoverable; the console asks for at least 8 characters). Create it once:

\begin{sh}
nodwyn keychain init          # or: the app's Settings page
\end{sh}

Then store a key from a session. The input is masked, never logged, never
echoed:

\begin{txt}
you> :eval (credential-register "weather-key")
\end{txt}

The front-end asks for the key in a masked prompt, stores it, and returns a
handle such as \code{\#<cred:weather-key:\ldots>}.

What comes back is a \textbf{handle}. It prints as
\code{\#<cred:class:id>} and nothing in the language can turn it into the
key: not \code{display}, not \code{format}, not \code{json-emit}, not
string concatenation. \code{credential-register} is user-only: neither the
model nor the loop's strategy can put a key prompt in front of you.
\code{(credential "weather-key")} looks up the handle of a key already
stored, or returns \code{\#f}; it is refused to the model, while the loop's
strategy may use it (the plan gate looks up the judge's key that way).

\section{Bindings: where a key may go}

A handle alone opens nothing. \file{bindings.toml} says where each class of
key may be sent, and how:

\begin{toml}
# ~/.config/nodwyn/bindings.toml
[[bind]]
class = "weather-key"
dest = "api.weather.example"
positions = ["query:appid"]
methods = ["GET"]

[[bind]]
class = "openai-key"
dest = "api.openai.com"
positions = ["header:Authorization:Bearer"]
methods = ["POST"]
\end{toml}

A position is where the key is injected:

\begin{itemize}
  \item \code{header:NAME}, optionally with a scheme:
    \code{header:Authorization:Bearer};
  \item \code{query:NAME}, a query parameter;
  \item \code{body}.
\end{itemize}

A service that uses OAuth2 (Google, Microsoft, GitHub, Notion, or an MCP
server behind a login) gets an \code{oauth2} credential, authorized in the
browser from the command line, never from a session:

\begin{sh}
nodwyn credential add --oauth2 google --client-id ID --scopes drive.file
nodwyn credential add --oauth2 custom --client-id ID \
    --auth-url https://auth.example/authorize --token-url https://auth.example/token
\end{sh}

Its row may name the scopes it allows; a token granted any other is refused:

\begin{toml}
[[bind]]
class = "google"
dest = "www.googleapis.com"
positions = ["header:Authorization:Bearer"]
methods = ["GET", "POST"]
scopes = ["https://www.googleapis.com/auth/drive.file"]
\end{toml}

The keychain refreshes the access token just before a use the table
allowed and the session logs \code{cred.refresh}. If the provider refuses
the refresh, every use is denied until \code{nodwyn credential login
CLASS}; \code{nodwyn credential revoke CLASS} revokes it at the provider and
deletes it here.

With a \code{google} credential stored, sessions offer the model
\code{google-doc-upload}, which sends a document to Google Drive as a new
Google Doc (Chapter~\ref{ch:workspaces} has the details). It needs only
\code{POST} in that row, and only \code{drive.file}.

The binding table is checked on every
request, before anything is sent; an unbound combination is never sent.
Only you edit it: the file is policy, and \code{policy-write} is user-only
and shows you a diff first.

\section{\texttt{http-request}}

\begin{scm}
(http-request cred url method position headers body)
;=> (("status" . 200) ("body" . "..."))
\end{scm}

The key is injected at \code{position} on the way out, if the binding table
allows exactly (class, host, method, position). Handles cannot come from
JSON, so the model can never call \code{http-request} directly. You reach
it through a tool you register:

\begin{scm}
;; Look the key up once, while init.scm runs with your origin. The tool
;; closes over the handle; the model's calls could not look it up.
(define weather-key (credential "weather-key"))

(register-tool "weather"
  (lambda (city)
    (if (not weather-key)
        (error "no weather-key was stored when this session started")
        (let ((r (http-request weather-key
                               (string-append "https://api.weather.example/v1/now?q=" city)
                               "GET" "query:appid" '() "")))
          (if (= (assoc-ref r "status") 200)
              (assoc-ref r "body")
              (format "weather service answered ~a" (assoc-ref r "status"))))))
  '((doc . "Current weather for a city, as JSON from the weather service.")
    (params . (("city" "string" "a city name, e.g. Lisbon")))
    (egress . #t)))
\end{scm}

\begin{careful}
A registered tool's body runs with the origin of whoever called it. When
the model calls \code{weather}, the body runs as the model, and
\code{(credential \ldots)} is refused to model-originated code. That is why
the handle is looked up at the top of the file, at load time, and captured
by the closure (Chapter~\ref{ch:forms}). Holding a handle grants nothing by
itself: every use is still checked against the binding table.
\end{careful}

What comes back is tagged \code{external:http}, and it passes the same
redaction as every tool result before the model reads it.

\subsection{What the guard checks on the way out}

\begin{itemize}
  \item \textbf{The binding} for the handle: class, host, method, position.
  \item \textbf{L1}: the headers and body are scanned for anything that
    looks like a credential. A recognized key there is refused unless its
    class is bound for that host and method (then the binding says it may
    go). Command lines are scanned the same way. The URL is the
    destination: its host is checked against the bindings, not scanned.
  \item \textbf{L2}: content that came from your files or command output may
    not go to a destination no binding names.
\end{itemize}

By default a keyed request that no binding covers is denied outright. With
\code{policy.unbound = "confirm"} in \file{config.toml}, it asks instead,
and the answer offers \key{B}: save this exact binding and allow. The binding is then in the
file, where you can read it and remove it.

\section{Models are egress too}

A model with a \code{credential} is a keyed destination like any other, and
the same binding must allow its host. Every request is checked and logged
as \code{egress.check}. When the app sets up a model for you, it writes
exactly that binding: the server's host, \code{POST}, a bearer
\code{Authorization} header.

An attached image is egress of a kind no binding governs: its pixels cannot
be redacted. It goes to the control model's host only, and only after a
confirm that names that host (Chapter~\ref{ch:workspaces} shows it); a
model you route to at another host gets a line saying the image is not for
it.

\section{MCP servers}

MCP servers offer tools. \nodwyn{} treats a server as a destination it
does not trust. In the console, \textbf{Connections \(\rightarrow\)
Tools} adds one, remote (HTTPS only; a key goes to the keychain, bound to
that host only) or contained on this computer (below), lists its tools for
review, and records what you decided in \file{mcp.toml}:

\begin{itemize}
  \item \textbf{A pinned manifest.} You approve the exact tool list: every
    name, description and schema is hashed. If the server later changes
    any of it, new tasks withhold its tools and say why, until you review
    again.
  \item \textbf{Declared effects.} For each tool you say whether it
    \emph{reads} or \emph{changes things}. A tool with no declaration is
    treated as both: it asks every call and offers no standing answer.
  \item \textbf{Nothing on by default.} A new server offers nothing, and a
    new tool is off until you turn it on.
\end{itemize}

The model sees an MCP tool as \code{server\_\_tool} (model APIs refuse
slashes in tool names); the console shows \code{server/tool}. Arguments are
egress to the server's host, and every one is a payload: L1 refuses a
recognized key in any of them, however deeply nested, unless its class is
bound for that host, and L2 lets content from your files or command output
go there because you approved the server, with a key or without. Results
are ingress: redacted, tagged \code{external:mcp}, and never instructions.

\begin{tip}
The confirm for an MCP call is not really ``is this argument safe?''. It
is ``do I want this third party to see this?''. The prompt names the
server, the host and what is being sent, so you can answer that question.
\end{tip}

\subsection*{A server on this computer}

A local server never runs as a process of yours. \nodwyn{} runs it in a
container, with podman or docker (\code{[mcp] runtime} in
\file{config.toml}), from an image pinned by digest; a tag is refused. You
choose what it can reach:

\begin{itemize}
  \item \textbf{Network.} \code{network = []}: none at all. A list of hosts:
    its only way out is \nodwyn{}'s egress proxy, which lets through
    exactly those hosts (ports 443 and 80, or \code{host:port}) and records
    every attempt as \code{mcp.egress}; the task shows the refused ones.
  \item \textbf{Files.} \code{mounts}: workspace folders it sees under
    \code{/workspace}, read-only unless \code{:write}. Never the workspace
    itself, never \code{.nodwyn/}.
  \item \textbf{A key.} \code{position = "env:NAME"}: the variable it
    reads its key from. The key goes in at start only if its class is bound
    for every listed host at that position; with no network it is refused.
\end{itemize}

A task never downloads code: \code{nodwyn mcp pull NAME}, or Pull in
Connections, fetches the image by its digest. Everything else is fixed: no
capability, a read-only root, none of your environment, and the container
is removed when the task ends. A confirm names ``the contained server NAME
(network: \ldots)''. An HTTP client in the server that ignores
\code{HTTPS\_PROXY} reaches nothing.
